% Generated draft: requires mathematical and source review.
\documentclass[11pt]{article}
\usepackage[margin=25mm]{geometry}
\usepackage{fontspec}
\setmainfont{Latin Modern Roman}
\usepackage{amsmath,amssymb,mathtools,unicode-math}
\setmathfont{Latin Modern Math}
\usepackage{enumitem,needspace,xurl,hyperref}
\setlength{\parindent}{0pt}
\setlength{\parskip}{5pt}
\setlength{\emergencystretch}{3em}
\newcommand{\N}{\mathbb{N}}
\newcommand{\Z}{\mathbb{Z}}
\newcommand{\Q}{\mathbb{Q}}
\newcommand{\R}{\mathbb{R}}
\newcommand{\C}{\mathbb{C}}
\newcommand{\F}{\mathbb{F}}
\newcommand{\A}{\mathbb{A}}
\newcommand{\PP}{\mathbb{P}}
\newcommand{\E}{\mathbb{E}}
\newcommand{\eps}{\varepsilon}
\newcommand{\dd}{\,\mathrm{d}}
\begin{document}
\section*{UnsolvedMath ID: 30000131}
\textbf{Source problem code:} catalogue entry
\subsection{Definitions and mathematical statement}
Let $f_1,f_2:M\to N$ be maps with the common domain and target intended in the report. For any pair of representatives define its coincidence set by
\[
 C(f'_1,f'_2)=\{x\in M:f'_1(x)=f'_2(x)\}.
\]
Write $\pi_0(C)$ for its component set, with the usual path-component convention, and $\#\pi_0(\varnothing)=0$. The required homotopy invariant is
\[
 M\!CC(f_1,f_2)=
 \min\left\{\#\pi_0(C(f'_1,f'_2)):
            f'_1\simeq f_1,\ f'_2\simeq f_2\right\}.
\]
The homotopies of the two maps are independent, within their respective original homotopy classes and with the same $M,N$. The question is to determine this minimum from the homotopy data, rather than merely restating its definition. In a setting permitting infinite component counts, one can record infinity if no finite count is possible; the source's intended finiteness hypotheses must be checked.

The extract does not give dimensions or compactness assumptions on the source and target manifolds, the allowed category of representatives and homotopies, or transversality conventions. In a smooth transverse representative pair, the coincidence set is a submanifold and connected and path components agree. Without such regularity, those component conventions need not agree, so no unprovided genericity restriction is inserted in the minimization.

The literature assessment checked on 21 August 2026 reports that normal-bordism and Nielsen invariants yield sharp formulas in several dimension ranges and target classes. It also reports examples in which the lower bound is not sharp and secondary obstructions persist, so a general formula is not known in that assessment. Counting coincidence components is different from counting coincidence points or simply declaring a Nielsen lower bound to be the exact minimum. These are dated literature-triage claims requiring independent verification.
\subsection{Short English statement}
Determine the smallest number of coincidence components obtainable by independently deforming each map within its homotopy class. Known Nielsen and bordism lower bounds are sharp in some cases but do not supply a general formula.
\subsection{Sources}
\begin{itemize}
\item \emph{Topologie} (2005), Oberwolfach Reports, DOI 10.4171/owr/2004/44, pp. 2311--2348, \url{https://doi.org/10.4171/owr/2004/44}; proposed year supplied as 2004.
\item Literature assessment checked 21 August 2026: Koschorke, ``Nielsen coincidence theory,'' \url{https://arxiv.org/abs/math/0408044}; Koschorke, ``Nielsen numbers in coincidence theory,'' \url{https://arxiv.org/abs/1305.1669}; original OWR report, \url{https://doi.org/10.4171/owr/2004/44}.
\item Local UnsolvedMath registry, permanent ID $30000131$, exact source code \texttt{catalogue entry}, created 25 August 2026 and updated 16 September 2026. The explicit minimization statement is labelled corrected and verified with malformed homotopy symbols restored.
\end{itemize}
\end{document}
